%% ===========================================================================
%% Evaluating Multi-Agent Communication Topologies for Automated Code Review
%% Revised work-in-progress draft. TARGET: 5 pages max, IEEE conference format.
%%
%% Overleaf: New Project -> paste this file. Compiler: pdfLaTeX.
%% IEEEtran.cls ships with Overleaf (official template:
%% https://www.ieee.org/conferences/publishing/templates). Insert your own
%% architecture diagram where Fig. 1 is commented in Section IV-B.
%% ===========================================================================
\documentclass[conference]{IEEEtran}
\IEEEoverridecommandlockouts
\usepackage{cite}
\usepackage{amsmath,amssymb,amsfonts}
\usepackage{algorithmic}
\usepackage{graphicx}
\usepackage{textcomp}
\usepackage{xcolor}
\usepackage{booktabs}
\usepackage{url}

\begin{document}

\title{Evaluating Multi-Agent Communication Topologies\\
for Automated Code Review: A Work-in-Progress Study}

\author{\IEEEauthorblockN{En-Ping Su, Tong Wu, Shiting Huang, Mengshan Li}
\IEEEauthorblockA{\textit{Khoury College of Computer Sciences} \\
\textit{Northeastern University}, Vancouver, Canada}}

\maketitle

\begin{abstract}
Large Language Models (LLMs) are increasingly used for automated code review, and
recent multi-agent systems suggest that dividing work among specialized agents can
improve reasoning on complex tasks. It remains unclear, however, whether the
\emph{communication topology} among agents---not merely their presence---affects
review quality, cost, and efficiency. This work-in-progress paper asks how
communication topology affects the effectiveness and efficiency of automated code
review. We compare three approaches under identical conditions: an Agentless
single-LLM baseline, a Hierarchical Authority architecture, and a Decentralized
Peer Consensus architecture. We have implemented a modular platform that reviews
the same immutable pull-request snapshot with all three architectures through a
shared import, review, validation, storage, and evaluation pipeline backed by a
production LLM service, isolating topology as the only independent variable.
Evaluation spans three complementary datasets: Qodo PR-Review-Bench (objective
ground truth), SWE-PRBench (agreement with human reviewers), and an actively
developed full-stack data portal used as an industrial case study. Because the
industrial repository lacks authoritative ground truth, we introduce an
\emph{industrial verification} methodology that corroborates findings through
cross-architecture agreement, static-analysis agreement, independent LLM-judge
validation, and later-fix rate rather than precision and recall. Preliminary runs
show communication overhead and latency rising steeply from Agentless to
Hierarchical to Consensus, framing the accuracy/cost trade-off the full study will
quantify.
\end{abstract}

\begin{IEEEkeywords}
Large Language Models, Multi-Agent Systems, Automated Code Review, Software
Engineering, Evaluation Methodology.
\end{IEEEkeywords}

\section{Introduction}
Code review is a central quality-assurance practice: before code is merged,
reviewers examine pull requests (PRs) for functional defects, security problems,
maintainability issues, and architectural inconsistencies. It is also
time-consuming and hard to scale, especially for full-stack systems spanning
frontend, backend, databases, and deployment infrastructure. LLMs can summarize
changes, flag bugs, and recommend fixes, and multi-agent frameworks such as
ChatDev and MetaGPT suggest decomposing complex tasks among specialized
agents~\cite{chatdev,metagpt}.

Existing work, however, usually evaluates \emph{complete} agent systems rather than
isolating the role of communication topology. A centralized manager may behave
very differently from a decentralized group in which agents discuss findings
directly: centralized coordination may be faster and cheaper, while decentralized
discussion may catch more cross-component issues. We therefore study whether the
\emph{structure} of agent communication---holding the model, prompts, and roles
fixed---changes review outcomes:

\begin{quote}
\textbf{RQ:} How does the communication topology of a multi-agent system affect the
effectiveness and efficiency of automated code review?
\end{quote}
We refine this into three sub-questions:
\begin{itemize}
  \item \textbf{RQ1 (Effectiveness).} Does topology change defect-detection
  quality (precision, recall, localization) and agreement with human reviewers?
  \item \textbf{RQ2 (Efficiency).} How does topology affect latency, token cost,
  and coordination overhead (LLM calls, inter-agent messages)?
  \item \textbf{RQ3 (Complexity interaction).} Does the effect of topology depend
  on PR complexity, e.g., localized versus cross-component changes?
\end{itemize}

We compare Agentless, Hierarchical Authority, and Decentralized Peer Consensus.
Agentless is a strong single-LLM baseline, since simple non-agent pipelines can be
competitive on software-engineering tasks~\cite{agentless}; the two multi-agent
systems share specialist roles but differ in how agents communicate. Our
contributions are:
\begin{itemize}
  \item A controlled design and an \emph{implemented}, reproducible platform in
  which review architecture is the only independent variable.
  \item A three-dataset evaluation framework combining objective benchmark
  correctness, human-reviewer agreement, and an industrial case study.
  \item An \emph{industrial verification} methodology that corroborates findings
  without authoritative ground truth.
  \item Preliminary observations on the communication-cost trade-off across
  topologies.
\end{itemize}

\section{Related Work}
\subsection{LLM-Based Software Engineering}
LLMs have been applied to code generation, program repair, bug localization, and
test generation. SWE-bench benchmarks real GitHub issues and PRs and shows the
difficulty of resolving realistic problems with language models~\cite{swebench}.
Agentless challenges the assumption that complex agent systems are always
necessary, using a simple localize--repair--validate pipeline~\cite{agentless};
its strong performance makes it an essential baseline for judging whether
multi-agent communication adds measurable value.

\subsection{Multi-Agent Software Engineering}
Multi-agent systems assign roles to multiple LLM agents. ChatDev models
development as communication among design, coding, and testing
agents~\cite{chatdev}; MetaGPT adds structured workflows and role-based
collaboration to reduce inconsistency~\cite{metagpt}. Surveys argue the area is
promising but under-evaluated empirically~\cite{survey}, and other work finds that
agent systems fail through coordination overhead, cascading errors, and weak
verification~\cite{failures}. These failure modes are precisely properties of
communication structure, motivating our controlled comparison of topologies rather
than of whole systems.

\subsection{Evaluation Without Ground Truth}
Objective benchmarks require labeled defects, which real industrial repositories
rarely provide. Using a separate LLM as an impartial judge is an increasingly
common way to assess outputs at scale~\cite{llmjudge}. We adopt LLM-as-a-judge as
one \emph{corroboration} signal---not a source of ground truth---and combine it
with cross-architecture agreement and conventional static analysis to triangulate
the credibility of findings on a live repository.

\section{Research Design}
\subsection{Variables}
The independent variable is communication topology across three architectures:
(1)~\textbf{Agentless}: one LLM reviews the PR without agent communication;
(2)~\textbf{Hierarchical Authority}: a manager assigns work to specialists and
synthesizes their findings; (3)~\textbf{Decentralized Peer Consensus}: specialists
review independently, exchange findings, and vote to converge. Controlled
variables (model, prompts, roles, PR snapshot) are held fixed; dependent variables
capture effectiveness and efficiency (Table~\ref{tab:vars}).

\begin{table}[t]
\caption{Research Variables}
\label{tab:vars}
\centering
\begin{tabular}{@{}ll@{}}
\toprule
\textbf{Type} & \textbf{Description} \\
\midrule
Independent   & Communication topology \\
Controlled    & Model, prompts, roles, PR snapshot \\
Effectiveness & Precision, recall, F1, localization; corroboration \\
Efficiency    & Latency, token cost, LLM calls, message count \\
\bottomrule
\end{tabular}
\end{table}

\subsection{Hypotheses}
\textbf{H1:} Hierarchical Authority has lower latency and token cost than
Decentralized Peer Consensus (RQ2). \textbf{H2:} Decentralized Peer Consensus
achieves higher recall on complex or cross-component PRs (RQ1, RQ3).
\textbf{H3:} Agentless is competitive on simple, localized PRs (RQ1, RQ3).
\textbf{H4:} The effect of topology grows with PR complexity (RQ3).

\section{Platform and Methodology}
\subsection{Experimental Platform}
We implemented a modular platform that drives every architecture through the same
pipeline: PR import produces an immutable snapshot; the architecture under test
reviews that snapshot; a validation stage normalizes raw model output into
structured findings; results are stored; and an evaluation engine computes metrics.
A campaign runner imports each PR once and reviews the shared snapshot with all
three architectures, so topology is the only variable (\emph{fairness policy}). All
architectures use the same foundation model, prompts, and roles; reviews run
against Anthropic Claude on Amazon Bedrock.

\subsection{Architectures}
% ===== FIGURE: upload your diagram to Overleaf as the filename below =====
% (spans both columns; place the image at the project root or in a figures/ folder)
\begin{figure*}[t]
  \centering
  \includegraphics[width=\textwidth]{architectures-comparison.png}
  \caption{The three review architectures compared. (a)~\textbf{Agentless}:
  a single LLM reviews the pull request with no agent communication.
  (b)~\textbf{Hierarchical Authority}: a coordinator assigns the diff to
  specialist agents (e.g., UI, API, data) and aggregates their findings---a
  one-to-many, coordinator-mediated pattern with no peer communication.
  (c)~\textbf{Decentralized Peer Consensus}: dynamically selected specialists
  review independently and then discuss directly (many-to-many) to resolve
  conflicts and produce a consensus report. Solid arrows denote control/task
  flow; dashed arrows denote peer communication.}
  \label{fig:arch}
\end{figure*}
Figure~\ref{fig:arch} contrasts the three architectures. Agentless reviews the
diff in a single pass. Hierarchical Authority has a coordinator (manager) direct
frontend, backend, and database specialists, then deterministically merges their
findings. Decentralized Peer Consensus uses the same specialists, which first
review independently, then exchange findings and vote to reach a shared set. The
multi-agent architectures share specialist implementations and prompts, differing
only in coordination and message flow.

\subsection{Datasets}
Three datasets provide complementary evidence: \textbf{Qodo PR-Review-Bench}
(PRs with injected defects and explicit ground truth $\rightarrow$ precision,
recall, F1, localization); \textbf{SWE-PRBench} (real PRs with human review
comments, interpreted as agreement with reviewers rather than absolute
correctness); and an \textbf{industrial case study}---an actively developed
full-stack data portal (React frontend, backend APIs, database logic,
authentication, cloud deployment)---reviewed on real PRs for ecological validity
but without authoritative ground truth.

\subsection{Evaluation Metrics}
Metrics are read according to the evidence each dataset offers
(Table~\ref{tab:metrics}). For the industrial case study, correctness cannot be
measured directly, so we do \emph{not} report precision/recall/F1. Instead,
\textbf{industrial verification} corroborates findings via:
(i)~\emph{Cross-Architecture Agreement}---the fraction of an architecture's
findings independently found by another architecture on the same PR (matched on
file, line, category, and title similarity);
(ii)~\emph{Static-Analysis Agreement}---the fraction coinciding with issues from
conventional static-analysis tools;
(iii)~\emph{LLM-Judge Validation}---the fraction a separate, impartial LLM
classifies as valid given the diff;
(iv)~\emph{Later-Fix Rate}---the fraction whose location is modified in later
commits; and
(v)~\emph{Evidence Score}, a supporting heuristic over severity, confidence, and
finding volume that reflects review \emph{strength}, not correctness. Operational
metrics (latency, tokens, cost, LLM calls, messages) are collected identically
across all datasets.

Formally, let $F$ be an architecture's validated findings on a PR and $A_x$ the
findings of architecture $x$. Writing $m(f,S)=1$ when a finding $f$ matches some
element of set $S$ (same file, a small line window, and matching category or
title), the corroboration signals are
\begin{align}
\mathrm{CAA}_x &= \tfrac{1}{|A_x|}\!\sum_{f\in A_x}\!
   \mathbb{1}\!\left[\exists\,y\neq x:\, m(f,A_y)\right], \label{eq:caa}\\
\mathrm{SAA}   &= \tfrac{1}{|F|}\!\sum_{f\in F}\!
   \mathbb{1}\!\left[m(f,S_{\mathrm{static}})\right], \\
\mathrm{LJV}   &= \tfrac{1}{|F|}\bigl|\{f\in F:\mathrm{judge}(f)=\text{valid}\}\bigr|, \\
\mathrm{LFR}   &= \tfrac{1}{|F|}\bigl|\{f\in F:\mathrm{loc}(f)\ \text{changed later}\}\bigr|,
\end{align}
where $S_{\mathrm{static}}$ is the set of static-analysis issues and
$\mathrm{judge}(\cdot)$ is an independent LLM verdict. Each value lies in $[0,1]$
and is reported only when computable---for example, cross-architecture agreement
(\ref{eq:caa}) requires at least two architectures and a non-empty finding set.

\begin{table}[t]
\caption{Evaluation Metrics by Dataset}
\label{tab:metrics}
\centering
\small
\begin{tabular}{@{}lccc@{}}
\toprule
\textbf{Metric} & \textbf{Qodo} & \textbf{SWE-PRBench} & \textbf{RAP Portal} \\
\midrule
Precision / Recall / F1        & \checkmark & $\circ$ & --- \\
Localization Accuracy          & \checkmark & ---     & --- \\
Human Review Agreement         & ---        & \checkmark & --- \\
Cross-Architecture Agreement   & ---        & ---     & \checkmark \\
Static-Analysis Agreement      & ---        & ---     & \checkmark \\
LLM-Judge Validation           & ---        & ---     & \checkmark \\
Later-Fix Rate                 & ---        & ---     & \checkmark \\
Evidence Score (heuristic)     & ---        & ---     & \checkmark \\
Latency / Cost / Tokens        & \checkmark & \checkmark & \checkmark \\
LLM Calls / Message Count      & \checkmark & \checkmark & \checkmark \\
\bottomrule
\end{tabular}
\\[2pt]
\raggedright
\footnotesize \checkmark: reported; $\circ$: interpreted as human agreement;
---: not applicable.
\end{table}

\subsection{System Implementation}
The platform is organized as a set of independent stages connected by explicit
data contracts, which keeps the experiment reproducible and makes each stage
separately testable. \emph{Import} converts a PR into an immutable snapshot keyed
by a content-derived identifier, so re-running an experiment reuses the exact same
input. \emph{Review} invokes the architecture under test, which returns raw model
output plus execution telemetry (latency, tokens, cost, LLM calls, messages).
\emph{Validation} parses the raw output against a fixed schema and normalizes it
into structured findings without inventing data, so malformed output degrades
gracefully rather than corrupting metrics. \emph{Storage} persists raw and
validated results and individual findings; \emph{evaluation} computes the metrics
of Section~IV-D; and \emph{export} emits per-experiment and comparison rows as CSV
and JSON for analysis.

Two properties support internal validity. First, the evaluation engine is
deterministic and side-effect free: it reads only stored artifacts and never calls
the model, so metrics are reproducible from saved results. Second, the campaign
runner maintains a manifest of every (PR, architecture, run) with status,
attempts, and reproducibility metadata (model, prompt, and workflow versions),
enabling resume-after-failure and auditing; a retry policy re-attempts only
transient infrastructure errors, while validation and parsing errors are terminal.
The industrial-verification signals are additive: they extend the research-evidence
metrics without altering the benchmark correctness path, so the Qodo and
SWE-PRBench pipelines are unaffected by the case-study methodology.

\section{Proposed Experiments}
\subsection{Campaign Design}
The primary experiment is a full campaign in which every PR is reviewed by all
three architectures over the shared snapshot. Because the pipeline reuses one
import per PR, all architectures see byte-identical input, so any outcome
difference is attributable to topology. Each run records reproducibility metadata
(model version, prompt version, workflow version) and a deterministic manifest so
runs can be resumed and audited; a retry policy re-attempts only transient
infrastructure failures, while validation and parsing errors are terminal.

\subsection{Datasets and Sample Sizes}
The planned campaign targets 30--50 PRs from Qodo, 20--30 from SWE-PRBench, and
10--20 from the industrial portal, each reviewed under all three architectures
(e.g., 50 PRs $\times$ 3 architectures $=$ 150 experiments). PRs are stratified by
type---frontend-only, backend-only, database-only, and cross-component---so that
RQ3 can be tested against complexity. Benchmark defect categories include missing
authorization checks, incorrect API contracts, frontend state bugs, query
inefficiencies, missing validation, error-handling problems, and cross-component
integration failures.

\subsection{Procedure and Analysis}
Each PR flows through import $\rightarrow$ review ($\times 3$) $\rightarrow$
validation $\rightarrow$ storage $\rightarrow$ evaluation $\rightarrow$ export.
Effectiveness is measured objectively on Qodo (precision, recall, F1,
localization), as human agreement on SWE-PRBench, and by corroboration on the
industrial portal. Efficiency is measured everywhere (latency, tokens, LLM calls,
messages).

Because the same three architectures review every PR, observations are paired,
and we will compare architectures with paired non-parametric tests (Wilcoxon
signed-rank for two-way contrasts, Friedman across all three) rather than assume
normality, reporting effect sizes alongside $p$-values and correcting for multiple
comparisons. To address RQ3 we stratify PRs by type (frontend-only, backend-only,
database-only, cross-component) and test for an interaction between topology and
complexity; H2 predicts that Consensus gains most on cross-component PRs. Cost is
reported both in absolute terms and as a cost-per-true-positive ratio, so that any
accuracy gain from richer coordination is weighed against its overhead. All
prompts and the model are frozen across architectures to preserve internal
validity, and every run stores its reproducibility metadata so the analysis can be
regenerated from saved artifacts.

\section{Results So Far}
The platform is implemented end-to-end and validated against the live service:
all three architectures run through the shared pipeline, the Qodo ground-truth
evaluation and the campaign runner are complete, and industrial findings are pulled
from real PRs with cross-architecture agreement and LLM-judge validation computed
automatically. During bring-up we found and fixed a defect in which the multi-agent
specialists returned no findings because their prompts omitted the expected
structured-output contract; once corrected, both multi-agent architectures produce
findings and recover seeded defects---evidence of the value of an observable
pipeline for methodology validation.

We report early, small-sample runs meant to validate the pipeline, not to draw
conclusions; statistical analysis awaits the full campaign. The most robust signal
is structural: communication overhead and latency grow sharply with topology
richness (Table~\ref{tab:comm}), consistent with H1. On a minimal seeded-defect
sample, all three architectures recovered the known defects (recall $=1.0$), while
macro precision decreased as fan-out grew (Table~\ref{tab:quality}) because the
multi-agent architectures surfaced more candidate findings, raising false
positives on these small PRs. This is the accuracy/cost tension the full study
must resolve: richer coordination need not improve precision, and can dilute it,
even as it multiplies cost. These figures are illustrative only, given the very
small current sample, and should not be read as final results.

\begin{table}[t]
\caption{Preliminary Benchmark Review Quality, Macro-Averaged Over the Current
Small Sample (illustrative; not statistically significant)}
\label{tab:quality}
\centering
\begin{tabular}{@{}lccc@{}}
\toprule
\textbf{Architecture} & \textbf{Precision} & \textbf{Recall} & \textbf{F1} \\
\midrule
Agentless    & 0.38 & 1.00 & 0.53 \\
Hierarchical & 0.20 & 1.00 & 0.33 \\
Consensus    & 0.31 & 1.00 & 0.44 \\
\bottomrule
\end{tabular}
\end{table}

\begin{table}[t]
\caption{Preliminary Communication and Cost per Architecture on a Representative
Industrial PR (single review; illustrative)}
\label{tab:comm}
\centering
\begin{tabular}{@{}lcccc@{}}
\toprule
\textbf{Architecture} & \textbf{LLM Calls} & \textbf{Messages} &
\textbf{Latency} & \textbf{Cost} \\
\midrule
Agentless    & 1 & 1  & $\sim$4\,s   & \$0.006 \\
Hierarchical & 3 & 8  & $\sim$55\,s  & \$0.052 \\
Consensus    & 9 & 22 & $\sim$150\,s & \$0.146 \\
\bottomrule
\end{tabular}
\end{table}

Next steps are to scale samples to the planned sizes, run the full campaign across
all three datasets, complete the static-analysis and later-fix producers for the
industrial case study, and analyze outcomes by architecture and PR type.

\section{Threats to Validity}
\textbf{Construct validity.} The industrial case study lacks authoritative ground
truth, so our signals corroborate rather than prove correctness, and the LLM judge
may share biases with the reviewer models. We mitigate this by triangulating four
independent signals of different kinds (peer agreement, static analysis, an
independent judge, and later fixes), by treating the Evidence Score only as a
supporting heuristic, and by never reporting benchmark correctness for that
dataset. On SWE-PRBench, matching metrics measure agreement with reviewers rather
than absolute correctness, and we label them as such.

\textbf{Internal validity.} Differences in outcome could stem from confounds other
than topology. We hold the model, prompts, and specialist roles fixed and review a
byte-identical immutable snapshot per PR across all three architectures, so
topology is the only manipulated variable; the deterministic evaluation engine and
stored reproducibility metadata let us regenerate every metric from saved results.

\textbf{External validity.} Synthetic defects may not fully represent production
defects, and a single primary industrial codebase limits generalizability. The
three-dataset design partly offsets this by spanning injected defects, real human
review, and a live repository; future work will add further repositories and
multiple LLM families.

\textbf{Conclusion validity.} The current sample is small and used only for
pipeline validation, so the reported observations are not statistically
significant. The full campaign uses the planned sample sizes and paired
non-parametric tests with effect sizes and multiple-comparison correction before
any claim is drawn.

\section{Conclusion}
This work-in-progress paper studies how multi-agent communication topology affects
automated code review. We built a controlled, reproducible platform that evaluates
Agentless, Hierarchical Authority, and Decentralized Peer Consensus under identical
conditions across three complementary datasets, and introduced an industrial
verification methodology for the realistic setting without authoritative ground
truth. Preliminary runs show communication overhead and latency rising steeply from
Agentless to Hierarchical to Consensus, framing the central accuracy/cost trade-off
the full study will quantify. The expected contribution is practical guidance for
designing LLM-based code-review systems that balance accuracy, cost, and
coordination complexity.

\begin{thebibliography}{00}
\bibitem{chatdev} C. Qian \emph{et al.}, ``ChatDev: Communicative agents for
software development,'' in \emph{Proc. ACL}, 2024.
\bibitem{metagpt} S. Hong \emph{et al.}, ``MetaGPT: Meta programming for a
multi-agent collaborative framework,'' in \emph{Proc. ICLR}, 2024.
\bibitem{agentless} C. S. Xia, Y. Deng, S. Dunn, and L. Zhang, ``Agentless:
Demystifying LLM-based software engineering agents,'' \emph{arXiv:2407.01489},
2024.
\bibitem{swebench} C. E. Jimenez \emph{et al.}, ``SWE-bench: Can language models
resolve real-world GitHub issues?'' in \emph{Proc. ICLR}, 2024.
\bibitem{survey} X. Zhang \emph{et al.}, ``LLM-based multi-agent systems for
software engineering: Literature review, vision and the road ahead,''
\emph{arXiv preprint}, 2024.
\bibitem{failures} Y. Li \emph{et al.}, ``Why do multi-agent LLM systems fail?''
\emph{arXiv preprint}, 2025.
\bibitem{llmjudge} L. Zheng \emph{et al.}, ``Judging LLM-as-a-judge with MT-Bench
and Chatbot Arena,'' in \emph{Proc. NeurIPS}, 2023.
\end{thebibliography}

\end{document}
